
\subsection{custom numbering}
\verb|\renewcommand{\theequation}{S.\arabic{equation}}| um S.1 Label zu haben
\begingroup
\renewcommand{\theequation}{S.\arabic{equation}}
\newcommand{\groesse}[1]{\hspace{1cm}\makebox[4cm][l]{#1}}
Die Beschreibung/Größe wurde in der Zeile in eine Box gesetzt, die um 1cm eingerückt ist. In der letzten Zeile ohne Box.
\begin{flalign}
      &\groesse{Arithmetisches Mittel}   & \bar{x} &= \mathrm{Mean}(\{x_i\}_{i=1}^n)=\frac{1}{n} \sum_{i=1}^{n} x_i \label{eq:arithmetisches_mittel} \\
      &\text{Arithmetisches Mittel}   & \bar{x} &= \mathrm{Mean}(\{x_i\}_{i=1}^n)=\frac{1}{n} \sum_{i=1}^{n} x_i 
\end{flalign}
\endgroup
Nach endgroup ist S.1 beendet, aber der counter muss zurückgesetzt werden:
\verb|\setcounter{equation}{0}|
\setcounter{equation}{0}
\begin{equation}
  Normales Label
\end{equation}

\subsection{siunitx}
\[\num[exponent-mode=scientific]{12300} \quad \num[exponent-mode=scientific]{12300(5)} \quad \num[exponent-mode=scientific,exponent-product=\times]{12300}\]
\begin{gather}
\num[round-mode = uncertainty,round-precision = 2]{12.345(678)}\\
\num[round-mode = uncertainty,round-precision = 2]{12.345(123)}\\
\num[round-mode = uncertainty,round-precision = 1]{12.345}\\
\num[round-precision = 2]{12.345}\\
\end{gather}

\begin{gather}
\num{1.2(3)(4)} \\
\sisetup{uncertainty-descriptors = {sys, stat}}
\num{1.2(3)(4)} \\
\num[uncertainty-descriptor-mode = subscript]{1.2(3)(4)}
\end{gather}

\subsection{irgendwelche Gleichungen}
\begin{equation}
\begin{array}{c@{\hspace{2cm}}c}
  U_\text{SP-SP}=U_\text{max}-U_\text{min} & \Delta U_\text{SP-SP}=\sqrt{2}\Delta U_i=\dfrac{\sqrt{2}}{25}U_\text{div}
\end{array}
\label{eq:U_SP}
\end{equation}

\begin{gather}
  equation1 \\
  \Delta U_\text{DC, Offset}=\sqrt{\left(\frac{\Delta U_\text{max}}{2}\right)^2+\left(\frac{\Delta U_\text{min}}{2}\right)^2}=\frac{1}{25\sqrt{2}} U_\text{div}
\end{gather}
Hier wurden (1) und (2) per Hand geschrieben:
\begin{gather}
\begin{array}{c@{\hspace{3cm}}c}
(1)\quad s=\frac{1}{2}at^2+v_0t &v=at+v_0\Leftrightarrow t=\dfrac{v-v_0}{a}\quad (2)\nonumber\\
\end{array}\\
(2)\text{ in }(1):\quad v^2=v_0^2+2as
\label{eq:v^2}
\end{gather}

\begin{subequations}
\begin{equation}
  \operatorname{min}_{a,b,c} 
  \frac{1}{2}\mathbf{w}^{T}\mathbf{w} + C \sum_{i=1}^{l}\xi_{i} 
\end{equation}    
\begin{equation}
  y_{i}\left(\mathbf{w}^{T}\phi(x_{i})+b\right)
\end{equation}
\end{subequations}


\subsection{aligning with align}
Doppeltes \verb|&&| erhöht den Abstand auf das maximal mögliche - verteilt also die Inhalte. Einfaches \verb|<rechtsbündig>&<linksbündig>| sagt: Richte alles links vom \verb|&| rechtsbündig und alles rechts vom \verb|&| linksbündig an dem Zeichen aus. Damit ist die erste Spalte, bestehend aus rechts- und linksbündig beendet. Das nächste \verb|&| eröffnet eine neue Spalte.

\verb|&&&| zwischen den Spalten, \verb|&| natürlich für Ausrichtung der Gleichheitszeichen.
\begin{align}
  A&=B&&& C&=D\\
  1&=2&&& 3&=4
\end{align}
Hier eröffnet jetzt das erste \verb|&| des Tripels die neue Spalte, die dann mit \verb|&&| gleichmäßig verteilt wird.

\verb|&| zwischen den Spalten
\begin{align}
  A&=B&C&=D\\
  1&=2&3&=4
\end{align}
Hier öffnet das \verb|&| die neue Spalte und macht ein bisschen default abstand. 
\begin{align}
  A&B
\end{align}

\paragraph{align} \verb|&&| zwischen allen Spalten
\begin{align}
  A&&B&&C
\end{align}
\verb|&&| und \verb|&&&&| vor letzter Spalte
\begin{align}
  A&&B&&&&C
\end{align}
\verb|&&| und \verb|&&&| vor letzter Spalte 
\begin{align}
  A&&B&&&C
\end{align}
mit \verb|&| zwischen Spalten
\begin{align}
  A&B&C
\end{align}
mit \verb|&| und \verb|&&&| vor letzter Spalte
\begin{align}
  A&B&&&C
\end{align}


\paragraph{flalign} mit \verb|&&| zwischen allen Spalten
\begin{flalign}
  A&&B&&C
\end{flalign}
mit \verb|&|
\begin{flalign}
  A&B&C
\end{flalign}
\verb|&&| und \verb|&|
\begin{flalign}
  A&&B&C
\end{flalign}

\subsection{Mathtexting}
\subsubsection{overset und overbracket}
\[\overset{def}{=}\]
How to use $\overbracket{abcd}^{e}$ \clink{https://tex.stackexchange.com/questions/132526/overbrace-and-underbrace-with-square-bracket}
\[\overbracket{1}_{indexunten}^{indexoben}\]

Eigener Custom command für Gleichheitszeichen. 
\newcommand{\mytexteq}[2]{%
  \mathrel{%
    \overset{\mbox{\normalfont\tiny\sffamily #1}}{%
      \underset{\mbox{\normalfont\tiny\sffamily #2}}{=}%
    }%
  }%
}
\[1\cdot 2\mytexteq{\(x\to\infty\)}{\(\epsilon\uparrow0\)}B\cdot A\]


\subsection{Boldsymbol in Heading}
\subsubsection[\texorpdfstring{\(K_\alpha\) lines}{K-alpha lines}]{\(\boldsymbol{K_\alpha}\) lines}